% HISTORICAL assembly fragment; the current integrated and checked text is
% parts/bmodels/bmodels.tex, which the master actually includes.
% AI-produced by GPT-6 Astra (OpenAI, ChatGPT Pro mode), 2026-11-19.
% Fragment for the root paper. Environments and natbib are supplied there.

\subsection{B.3. Finite figures and event structures: the missing distributivity hypothesis}

\begin{definition}[The two distinct irreducibility conditions]\label{SA21-BM-def-prime}
Let \((\mathcal F,\leq)\) be a nonempty partially ordered set with least
element \(0\) and suprema of bounded families. An element \(m\ne0\) is
\emph{completely join-irreducible} if
\[
 m=\bigvee A\quad\Longrightarrow\quad m\in A
\]
whenever the supremum exists. It is \emph{prime for bounded joins} if
\[
 m\leq\bigvee A\quad\Longrightarrow\quad
 (\exists a\in A)\ m\leq a
\]
for every bounded family. The definition in 156-4, scan 54, is the
first condition, with equality on the left. Write \(M\) for these
irreducible figures and \(\delta(F)=\{m\in M:m\leq F\}\) for deployment.
Here a \emph{finite figure system} means that each \(\delta(F)\) is
finite, not that \(M\) or \(\mathcal F\) is finite. We retain C1--C3:
bounded suprema exist, \(F=\bigvee\delta(F)\), and two figures have
a common upper bound exactly when every pair of their multistrates
has a common upper bound.
\end{definition}

\begin{counterexample}[Refuted: R12 and the literal form of L1]\label{SA21-BM-r12}
The literal finite-figure axioms C1--C3 do not imply representation
by finite downward-closed compatibility cliques. The failure persists
after adding C4, C4 bis, and the C5--C8 subdivision axioms, including
C8\(_0\), as stated in 156-4, scans 65--87.
\end{counterexample}

\begin{proof}
Take the six-element lattice
\[
 \mathcal F=\{0,d,a,b,c,1\},\qquad 0<d<a,b,c<1,
\]
where \(a,b,c\) are pairwise incomparable. All omitted comparisons
follow by reflexivity or transitivity. Distinct members of
\(\{a,b,c\}\) have join \(1\) and meet \(d\). All families have
suprema, so C1 holds. The equality-based definition gives
\[
 M=\{d,a,b,c\},\quad
 \delta(0)=\varnothing,\quad\delta(d)=\{d\},\quad
 \delta(u)=\{d,u\}\ (u=a,b,c),\quad\delta(1)=M.
\]
These sets generate their figures, proving C2. Every pair of figures
has upper bound \(1\); compatibility is universal and C3 holds.
The incidence poset \(M\) is connected, since \(d<a,b,c\).
Consequently C4 is vacuous and C4 bis is the identity decomposition
into one component.

Define subdivision by equality. The refinement induced by taking
a subfigure of a subdivision is then \(\leq\). C5, invariance of
disjointness under subdivision, is immediate. In C6, the subdivision
of \(G\leq F\) induced from \(F'=F\) is uniquely \(G'=G\).
Induction is the identity on the principal ideal of \(F\), so
preserves every bounded supremum and satisfies the multistrate
formula. For C7\(_0\), \(G'=G\) extends by \(F'=F\); the stronger
canonical choice is equally forced. A compatible family of
subdivisions in C8 consists of the original cover members, whose
given supremum is its own subdivision. Both zigzag conditions of
C8\(_0\) hold because equality in either subdivision position
forces equality throughout. Moderation holds since each figure
has only its identity subdivision. These verifications agree
with the observation in 156-4, scan 81, that C5--C7 are
tautological for equality subdivisions.

Nevertheless \(\{d,a,b\}\) is a finite downward-closed compatibility
clique in \(M\) and is not a deployment:
\[
 \delta(a\vee b)=\{d,a,b,c\}
 \ne\delta(a)\cup\delta(b)=\{d,a,b\}.
\]
The assertion in 156-4, scan 55, that deployment sends suprema
to unions is therefore false under its preceding hypotheses.
More intrinsically, this lattice is not the finite configuration
poset of any prime event structure. The ideal below a
configuration is the downset lattice of its causal poset and
is distributive, whereas here
\[
 a\wedge(b\vee c)=a,\qquad
 (a\wedge b)\vee(a\wedge c)=d.
\]
The program \texttt{bmodels\_checks.py} checks all families in
this specified lattice and the displayed axioms. This is a
counterexample to the C1--C8 package; it does not satisfy the
later requirement of a lieu in every open shadow.
\end{proof}

\begin{theorem}[Proved: exact corrected event-structure representation]\label{SA21-BM-event-iff}
For a finite-figure system satisfying C1--C3, the following are equivalent:
\begin{enumerate}
\item Every principal ideal \(\mathcal F_{\leq F}\) is distributive.
\item Every multistrate is prime for bounded joins.
\item Deployment is an isomorphism of ordered sets
\[
 \mathcal F\cong
 \{A\subseteq M:A\text{ is finite, downward closed, and an }R
   \text{-clique}\},
\]
where \(mRn\) means that \(m,n\) have a common upper bound.
\item \(\mathcal F\) is the finite configuration poset of a
binary-conflict prime event structure with finite causes.
\end{enumerate}
The event structure is recovered, uniquely up to a bijection
preserving causality and conflict, by
\[
 E=M,\qquad \leq_E=\leq|_M,\qquad
 m\mathbin{\#}n\ \Longleftrightarrow\ \neg(mRn).
\]
These conditions identify the whole representable subclass of
literal finite contrées. Adding C4 corresponds exactly to
requiring all conflict to occur within undirected causality
components. Adding C4 bis while retaining finite figures
requires that there be only finitely many nonempty such components.
\end{theorem}

\begin{proof}
C2 makes deployment order reflecting: if
\(\delta(G)\subseteq\delta(F)\), then
\(G=\bigvee\delta(G)\leq F\). Every principal ideal is therefore
finite, since it injects into \(\mathcal P(\delta(F))\).
C1 makes it a lattice; a meet is the supremum of its common
lower bounds.

Assume (1). For \(m\leq G\vee H\), distributivity in a common
bounding ideal gives
\[
 m=(m\wedge G)\vee(m\wedge H).
\]
Irreducibility gives \(m\leq G\) or \(m\leq H\). Induct on a
finite family; any bounded family takes values in a finite
ideal, so this also proves the arbitrary bounded version.
The empty family causes no exception since \(m\ne0\).
This is the local bounded form of
\citet[Chapter IX, \S4, Lemma 1 and Theorem 5, p.~139]{Birkhoff1948}.

Assume (2). Then
\(\delta(\bigvee_iF_i)=\bigcup_i\delta(F_i)\) for bounded
families. Deployments are finite downward-closed cliques.
Conversely let \(A\subseteq M\) be such a clique. Construct
its join successively. Once the join of \(a_1,\ldots,a_k\)
exists, its deployment is
\(\bigcup_{i\leq k}\delta(a_i)\). Every member of this union
is compatible with \(a_{k+1}\), by downward heredity of
compatibility. C3 makes the partial join compatible with
\(a_{k+1}\), so C1 supplies their join. Thus \(\bigvee A\)
exists, and primality and downward closure give
\[
 \delta(\bigvee A)=\bigcup_{a\in A}\delta(a)=A.
\]
Use \(0\) if \(A=\varnothing\). This proves (3).
Under (3), the ideal below \(F\) is the lattice of all downsets
of the finite poset \(\delta(F)\); meets and joins are
intersection and union, proving (1).

The relation \(R\) is reflexive, symmetric and downward
hereditary. Its complement \(\#\) is therefore symmetric,
irreflexive and upward hereditary. Each principal causal
ideal is \(\delta(m)\), hence finite and conflict free.
Thus (3) is exactly the finite configuration definition
of \citet[Definition 14, pp.~99--100]{NPW1981}, with finite
causes additionally imposed.

Conversely, for a binary event structure with finite causes,
the finite configurations are the finite downward-closed
conflict-free sets. A family bounded by a finite configuration
has union as its supremum. Each configuration is the union
of its principal causes \(\downarrow e\). These are precisely
the nonzero completely join-irreducible configurations:
membership of \(e\) proves their prime property, and
a configuration not equal to any of its principal causes
is their proper join. Compatibility is absence of cross
conflict, so C3 holds. Principal ideals are downset lattices
and are distributive. The maps \(e\mapsto\downarrow e\)
and \(F\mapsto\delta(F)\) make the reconstructions inverse
up to the stated isomorphisms. We identify objects and
isomorphisms here, not an unspecified category of morphisms.

Incidence components of \(M\) are causal components. C4
says precisely that events from different components
never conflict. Configurations then decompose uniquely
as component configurations with finite support.
If there are finitely many components, this is the
whole product, giving C4 bis. Conversely, choosing a
nonempty principal cause in each of infinitely many
components yields a tuple with infinite support, which
is not the trace of a finite configuration. When \(E\)
is empty the empty product is the one-element
configuration poset, as required.
\end{proof}

\begin{counterexample}[Refuted: omission of the word binary]\label{SA21-BM-binary}
Even after correcting R12 by local distributivity, the
corresponding class is the binary-conflict class, not all
prime event structures under the finite-consistency convention.
\end{counterexample}

\begin{proof}
\citet[Definition 1.3.4, pp.~335--336]{Winskel1987} allows
a finite consistency predicate \(\mathrm{Con}\).
Take three events \(e_0,e_1,e_2\), equality causality, and
\[
 \mathrm{Con}=\{A\subseteq\{e_0,e_1,e_2\}:|A|\leq2\}.
\]
Singleton consistency, downward closure, closure under
adding causes, and finite causes all hold. All pairs
are consistent but the triple is not. Hence the
configurations \(\{e_0,e_1\}\), \(\{e_2\}\) violate C3.
No binary relation represents this consistency
predicate, since consistency of all pairs would force
empty conflict. The program checks these axioms and
all eight undirected conflict relations on the three
events. The precise subclass is that of
\citet[\S3.4.6, p.~378]{Winskel1987}.
\end{proof}

\begin{proposition}[Proved: configuration domains and supports]\label{SA21-BM-domain}
For the corrected class, let \(\mathcal C_f(E)\) denote finite
configurations and \(\mathcal C(E)\) all downward-closed
conflict-free subsets. Then
\[
 \mathcal C(E)\cong\operatorname{Idl}(\mathcal C_f(E)),
\]
where ideals are nonempty directed downward-closed sets.
The compact elements are the finite configurations,
and the complete primes are \(\downarrow e\).
This is the coherent prime-algebraic domain correspondence
of \citet[Theorems 8--9, Definitions 17--18, p.~102]{NPW1981}.
For the whole-figure disjointness of 156-4 set
\[
 e\perp_E f\quad\Longleftrightarrow\quad
 \neg(e\#f)\text{ and }\downarrow e\cap\downarrow f=\varnothing.
\]
The support ortholattice of finite configurations is
canonically isomorphic to the closed double polars
of \((E,\perp_E)\), and to the support ortholattice
using all configurations. This concerns whole-figure
disjointness, not the later primitive internal disjointness.
If the empty configuration is retained as an orthogonality
vertex, it is universally self-orthogonal, and every closed
figure polar contains it. Removing this one radical vertex
gives the ordinary irreflexive orthoset convention and the
same support ortholattice.
\end{proposition}

\begin{proof}
Finite configurations contained in a configuration \(x\)
form an ideal: finite unions remain configurations.
Their union is \(x\), because every \(e\in x\) lies
in \(\downarrow e\subseteq x\), which is finite.
Conversely, the union of a directed ideal is downward
closed, and any two of its events lie together in
one member and so do not conflict. Each finite
subconfiguration of this union lies in a member
by finite directedness. These operations are
inverse order isomorphisms.

Directed suprema are unions. Finite configurations
are compact by the finite-witness argument.
An infinite configuration is the directed union
of its finite subconfigurations, so cannot be
compact. Membership of \(e\) proves complete
primality of \(\downarrow e\); decomposition
into these principals proves there are no
others. A pairwise compatible family of
configurations has conflict-free union, proving
coherence and checking the cited theorem's
hypotheses in this setting.

For configurations \(x,y\),
\[
 x\perp y
 \ \Longleftrightarrow\
 x\cup y\in\mathcal C(E),\ x\cap y=\varnothing
 \ \Longleftrightarrow\
 (\forall e\in x)(\forall f\in y)\ e\perp_Ef .
\]
The forward implication uses causal closure.
Conversely cross nonconflict gives a configuration
union; an event in \(x\cap y\) would imply
\(e\perp_Ee\), impossible.

For \(A\subseteq E\) write \(A^{\perp_E}\) for
its event polar and set
\(K_A=\{x\in\mathcal C_f(E):x\subseteq A\}\).
Event polars are downward closed, since
\(\perp_E\) is downward hereditary.
If \(A=A^{\perp_E\perp_E}\), every \(e\in A\)
occurs in \(\downarrow e\in K_A\), whence
\[
 K_A^\perp=K_{A^{\perp_E}}.
\]
For any configuration family \(\mathcal U\),
the pairwise criterion gives
\[
 \mathcal U^\perp=K_{(\bigcup\mathcal U)^{\perp_E}}.
\]
Thus \(A\mapsto K_A\), with inverse union,
is a bijection between closed event polars
and closed figure polars. It preserves
inclusion and orthocomplement and hence
all complete ortholattice operations.
The same argument with all configurations
proves the final assertion. In the convention retaining
the empty figure, \(K_\varnothing=\{\varnothing\}\)
is the bottom closed polar. Removing the empty figure
from every polar is an order-and-complement isomorphism
to the convention with bottom the empty family.
\end{proof}

\begin{proposition}[Proved: refinement is additional domain structure]\label{SA21-BM-refinement}
The event structure determines incidence, compatibility
and the preceding whole-figure supports, but not refinement,
subdivision, lieux, or primitive internal disjointness.
Given multistrate refinement satisfying At 4, its translation is
\[
 g\ll f\quad\Longleftrightarrow\quad
 (\forall e\in g)(\exists d\in f)\ e\ll d .
\]
Under At 13, subdivision translates as
\[
 g\preceq f\quad\Longleftrightarrow\quad
 g\ll f\text{ and }g^\perp=f^\perp,
\]
equivalently maximality for configuration inclusion among
refinements of \(f\).
\end{proposition}

\begin{proof}
The first formula is At 4 through deployment.
Equality of supports is equivalent to equality of
annihilator families by triple-polar equality.
At 13 gives the maximality equivalence of 156-6,
scans 84--85, proved in Part B.1.

For nonuniqueness, take two events \(p,q\) with
equality causality and conflict \(p\#q\).
Configurations are \(\varnothing,\{p\},\{q\}\).
On \(M=\{p,q\}\), let incidence be equality
and primitive internal disjointness empty.
Use equality refinement in the first system;
in the second also put \(p\ll q\).
Both are partial orders containing incidence.
The target's incidence ideal is a singleton,
so the least carrier exists and strong
refinement is the same relation. Mag 3 is
vacuous, and Mag 4 is vacuous on distinct
incident substrata. The first system has
lieux \(p,q\); the second only \(p\).
All open shadows are nonempty and any
pair of them contains nondisjoint lieux,
so Mag L2 holds. They have the same
event structure and different refinement
and lieu data. The program checks Mag
1--4, divisibility and Mag L2 in both.
\end{proof}
